\paragraph{Notation. } Let $[n]$ be a the set $\{1,\dots,n\}$ and $\Delta([n])$ denote the probability simplex over $[n]$. For the simplicity of notations, we let $\cV = [n]$. Let $S_n$ denote the permutation group of $\{1,\dots,n\}$. For any $\sigma \in S_n$, $\sigma^i$ means applying permutation $\sigma$ $i$-times. For any matrix $M \in \mathbb{R}^{n\times n}$, let $M(x,y), M(x,:), M(:,y)$ denote the entry on $x$-th row and $y$-th column, the $x$-th row, and the $y$-th column respectively. 

\section{Proof of \Cref{thm:log-growth}}
\label{sec:e-proof}
\begin{proof}
Let 
\begin{align*}
    \cR = \left\{r \in \R^{n \times n}: \sum_{s \in \cV} r(v,s) = 1, ~r(v,s) \geq 0,~ \forall v,s \in \cV \right\}.
\end{align*}
By \Cref{lem:normalizationconstant}, the original problem is equivalent to
\begin{align*}
    J': ~\sup_{r \in \cR} \inf_{q \in \cQ(p_0,\delta)} \sup_{w \in \cP(p_0,q)} \sum_{v,s \in \cV} w(v,s) \cdot \left(\log r(v,s) - \log p_0(s)\right),
\end{align*}
with $r(v,s) = {p_0(s)e(v,s)} \in \cR$.
Note that 
\begin{align*}
    \sum_{v,s \in \cV} w(v,s) \cdot (- \log p_0(s)) = \sum_{s \in \cV} p_0(s) \cdot (- \log p_0(s)) = H(p_0)
\end{align*}
where $H(p_0)$ is the entropy of $p_0$. 
By \Cref{prop:reform-problem}, the optimum value of $J'$ is equal to 
\begin{align*}
(1-\tfrac{\delta}{2})\log\left(1-\frac{\delta}{2}\right)
 + \tfrac{\delta}{2}\log\left(\frac{\delta}{2(n-1)}\right),
\end{align*}
achieved at
\begin{align*}
r^*(v,s)
=
\begin{cases}
1-\delta/2, & s=v,\\
\delta/(2n-2), & s\neq v,
\end{cases}
\end{align*}
and for any $q \in \cQ(p_0,\delta)$ written as $q = \sum_{i=1}^k \lambda_i \cdot q_i$ for $q_i = p_0 + \frac{\delta}{2}\cdot (\mathbf{e}_{v_i} - \mathbf{e}_{s_i}) $, the optimizer of the inner problem is given by
\begin{align*}
    \sum_{i=1}^k \lambda_i \cdot \left(\frac{\delta}{2} \cdot (\mathbf{e}_{v_i}\mathbf{e}_{s_i}^\top - \mathbf{e}_{s_i}\mathbf{e}_{s_i}^\top) + \diag(p_0)\right).
\end{align*}
Therefore, the optimum value of the original problem is equal to
\begin{align*}
(1-\tfrac{\delta}{2})\log\left(1-\frac{\delta}{2}\right)
 + \tfrac{\delta}{2}\log\left(\frac{\delta}{2(n-1)}\right) + H(p_0).
\end{align*}
In particular, this is achieved at
\begin{align*}
e^*(v,s)
=
\begin{cases}
\frac{1-\delta/2}{p_0(s)}, & s=v,\\
\frac{\delta}{2(n-1)p_0(s)}, & s\neq v,
\end{cases}
\end{align*}
and the optimizer of the inner problem is given in the same way. 
This completes the proof.
\end{proof}


\subsection{Supporting lemma}

\begin{lemma}[Diagonal Dominance]\label{lem:diagonal-dom}
    Let $\cV$ be a the set $\{1,\dots,n\}$. For any matrix $M_0 \in \mathbb{R}^{n\times n}$, there exists a permutation matrix $P \in \R^{n \times n}$ such that the following condition holds for any $v \in \cV$ and permutation $\sigma \in S_n$:
    \begin{align}\label{eq:diagdom}
        \sum_{i=0}^{\kappa_v-1} (M_0P)(\sigma^i(v),\sigma^i(v)) \geq \sum_{i=0}^{\kappa_v-1} (M_0P)(\sigma^i(v),\sigma^{i+1}(v)).
    \end{align}
    where $\kappa_v$ is the minimum positive integer such that $\sigma^{\kappa_v}(v) = v$. 
\end{lemma}


\begin{proof}
Let $\mathcal{P}$ be the set of all permutation matrices and
\begin{align*}
    P^* = \arg \sup_{P \in \mathcal{P}} \tr[M_0 P].
\end{align*}

Then we claim that $P^*$ is the permutation matrix s.t. $M = M_0P^*$ satisfying Eq.~\eqref{eq:diagdom}. 

Denote 
\begin{align*}
    C = \{v, \sigma(v),\dots,\sigma^{\kappa_v-1}(v)\}.
\end{align*}
Suppose for the sake of contradiction that there exists $\sigma(\cdot)$ and $v$ such that 
\begin{align}\label{eq:cycle_not_satisfied}
    \sum_{i=0}^{\kappa_v-1} M(\sigma^i(v),\sigma^i(v)) < \sum_{i=0}^{\kappa_v-1} M(\sigma^i(v),\sigma^{i+1}(v)).
\end{align}
Define permutation
\begin{align*}
    \bar \sigma(u) = \begin{cases}
        \sigma^{-1}(u),&~ u \in C,\\
        u, &~ u \notin C.
    \end{cases}
\end{align*}
and permutation matrix
\begin{align*}
    P_{\bar \sigma} = \begin{pmatrix}
        \mathbf{e}_{\bar \sigma(1)} & \cdots \mathbf{e}_{\bar \sigma(n)}
    \end{pmatrix}^\top
\end{align*}
Then we have
\begin{align*}
    M P_{\bar \sigma}  = &~ \sum_{i=1}^n M(:,i) \cdot \mathbf{e}_{\bar \sigma(i)}^\top\\
    = &~  \sum_{i \notin C} M(:,i) \cdot \mathbf{e}_{\bar \sigma(i)}^\top + \sum_{i \in C} M(:,i) \cdot \mathbf{e}_{\bar \sigma(i)}^\top\\ 
    = &~ M + \sum_{i \in C} M(:,i) \cdot (\mathbf{e}_{\bar \sigma(i)}^\top - \mathbf{e}_i^\top)\\
    = &~ M + \sum_{j=0}^{\kappa_v-1} M(:,\bar \sigma^j(v)) \cdot (\mathbf{e}_{\bar \sigma^{j+1}(v)}^\top - \mathbf{e}_{\bar \sigma^j(v)}^\top ) .
\end{align*}
Notice
\begin{align*}
    \tr[P_{\bar \sigma} M] = &~ \tr[M] + \tr[\sum_{j=0}^{\kappa_v-1} M(:,\bar \sigma^j(v)) \cdot (\mathbf{e}_{\bar \sigma^{j+1}(v)}^\top - \mathbf{e}_{\bar \sigma^j(v)}^\top )]\\
    = &~ \tr[M] + \tr[\sum_{j=0}^{\kappa_v-1} M(:,\bar \sigma^j(v)) \cdot \mathbf{e}_{\bar \sigma^{j+1}(v)}^\top ] - \tr[\sum_{j=0}^{\kappa_v-1} M(:,\bar \sigma^j(v)) \cdot \mathbf{e}_{\bar \sigma^j(v)}^\top ]\\
    = &~ \tr[M] + \sum_{j=0}^{\kappa_v-1} M(\bar \sigma^{j+1}(v), \bar \sigma^j(v)) - \sum_{j=0}^{\kappa_v-1}  M(\bar \sigma^j(v),\bar \sigma^j(v))\\
    = &~ \tr[M] + \sum_{j=0}^{\kappa_v-1} M(\sigma^{j}(v), \sigma^{j+1}(v)) - \sum_{j=0}^{\kappa_v-1}  M(\sigma^j(v),\sigma^j(v))\\
    > &~ \tr[M]
\end{align*}
where the last inequality follows from Eq.~\eqref{eq:cycle_not_satisfied}.

But $\mathcal{P}$ is a group so $P^* P_{\bar \sigma} \in \mathcal{P}$. This contradicts with the definition of $P^*$.
\end{proof}

\begin{lemma}[Optimizer of Inner Problem]\label{lem:inner}
Let $\cV$ be the set $\{1,\dots,n\}$ and $p_0 \in \Delta(\cV)$ be a distribution over $\cV$ such that $\inf_{v \in \cV}p_0(v) > \delta$.  
Let $M: \cV \times \cV \mapsto \R$ be a matrix that satisfies
\begin{align}
    \sum_{i=0}^{\kappa_v-1} M(\sigma^i(v),\sigma^i(v)) 
     \geq  \sum_{i=0}^{\kappa_v-1} M(\sigma^i(v),\sigma^{i+1}(v))
    \label{eq:M-cycle-cond}
\end{align}
for any $v \in \cV$, $\sigma \in S_n$, where $\kappa_v = \inf\{k\ge1:\sigma^{k}(v) = v\}$. 
Fix distinct $a,b\in\cV$ and define
\begin{align*}
    q = p_0 + \frac{\delta}{2}\cdot (\mathbf{e}_a - \mathbf{e}_b).
\end{align*}
Consider the optimization problem
\begin{align*}
    \sup_{w \in \cP(p_0,q)} J(w), 
    \qquad 
    J(w) := \sum_{v,s \in \cV} w(v,s)   M(v,s),
\end{align*}
where
\begin{align*}
    \cP(p_0,q) 
    = \left\{w:\cV\times\cV\to\R_+  \Big|  
        \sum_{s\in\cV} w(v,s) = q(v) \ \forall v\in\cV,~\sum_{v\in\cV} w(v,s) = p_0(s) \ \forall s\in\cV
      \right\}.
\end{align*}
Then there exists a permutation $\sigma \in S_n$ such that 
\begin{align}
    w^* 
    = \frac{\delta}{2} \sum_{i=0}^{\kappa-2} 
        \left(\mathbf{e}_{\sigma^{i}(a)}\mathbf{e}_{\sigma^{i+1}(a)}^\top 
              - \mathbf{e}_{\sigma^{i+1}(a)}\mathbf{e}_{\sigma^{i+1}(a)}^\top\right) 
      + \diag(p_0)
      \label{eq:p-star-form}
\end{align}
is an optimizer of $J$, where $\kappa$ is the minimum positive integer such that $\sigma^{\kappa}(a) = a$.
\end{lemma}

\begin{proof} First, we establish existence of an optimizer. Define $w^{\mathrm{feas}}$ by
\begin{align*}
    w^{\mathrm{feas}}(a,b) = &~ \frac{\delta}{2},\\
    w^{\mathrm{feas}}(v,v) = &~ 
    \begin{cases}
        p_0(b) - \frac{\delta}{2}, & v=b,\\
        p_0(v), & v\neq b,
    \end{cases}\\
    w^{\mathrm{feas}}(v,s) = &~ 0, \quad (v,s)\notin\{(a,b)\}\cup\{(v,v):v\in\cV\}.
\end{align*}
Because $\inf_v p_0(v)>\delta>\delta/2$, we have $w^{\mathrm{feas}}\ge0$.  
Its row sums satisfy
\begin{align*}
    \sum_s w^{\mathrm{feas}}(a,s) = &~ p_0(a)+\frac{\delta}{2} = q(a),\\
    \sum_s w^{\mathrm{feas}}(b,s) = &~ p_0(b)-\frac{\delta}{2} = q(b),\\
    \sum_s w^{\mathrm{feas}}(v,s) = &~ p_0(v) = q(v), \quad v\notin\{a,b\},
\end{align*}
and its column sums satisfy
\begin{align*}
    \sum_v w^{\mathrm{feas}}(v,b) = &~ p_0(b)-\frac{\delta}{2} + \frac{\delta}{2} = p_0(b),\\
    \sum_v w^{\mathrm{feas}}(v,s) = &~ p_0(s), \quad s\neq b.
\end{align*}
Thus $w^{\mathrm{feas}}\in\cP(p_0,q)$. The set $\cP(p_0,q)$ is a nonempty bounded polytope, and $J$ is linear, so there exists
\begin{align*}
    w^{(0)} \in \arg\sup_{p\in\cP(p_0,q)} J(w).
\end{align*}

\medskip
For a feasible $w$, let $G_w$ be the directed graph on vertex set $\cV$ with an edge
$(v,s)$ whenever $v\neq s$ and $w(v,s)>0$.
A directed cycle in $G_w$ is a $k$-tuple $(v_0,\dots,v_{k-1})$ of distinct vertices,
$k\ge2$, such that
\begin{align*}
    w(v_i,v_{i+1})>0,\quad i=0,\dots,k-1,
\end{align*}
with indices understood modulo $k$ (so $v_k=v_0$).

Next, we show that for an optimal distribution $\bar{w}$, the graph $G_{\bar{w}}$ has no directed cycles. Fix any feasible $w$ and such a cycle $(v_0,\dots,v_{k-1})$.
For $\varepsilon>0$ define $\tilde w$ by
\begin{align}
    \tilde w(v_i,v_{i+1}) = &~ w(v_i,v_{i+1}) - \varepsilon,
    &&i=0,\dots,k-1,
    \label{eq:cycle-transform-off}\\
    \tilde w(v_i,v_i)     = &~ w(v_i,v_i) + \varepsilon,
    &&i=0,\dots,k-1,
    \label{eq:cycle-transform-diag}\\
    \tilde w(x,y)         = &~ w(x,y),
    &&(x,y)\notin\{(v_i,v_{i+1}),(v_i,v_i):0\le i\le k-1\}.
    \notag
\end{align}
Choose
\begin{align*}
    0<\varepsilon\le\inf_{0\le i\le k-1} w(v_i,v_{i+1}),
\end{align*}
so that $\tilde w\ge0$.

\emph{Row sums.}
For each $i$,
\begin{align*}
    \sum_{s\in\cV} \tilde w(v_i,s)
    = &~ \left(w(v_i,v_{i+1})-\varepsilon\right)
       + \left(w(v_i,v_i)+\varepsilon\right)
       + \sum_{s\notin\{v_i,v_{i+1}\}} w(v_i,s)\\
    = &~ \sum_{s\in\cV} w(v_i,s).
\end{align*}
All other rows are unchanged, so
\begin{align}
    \sum_{s} \tilde w(v,s) = \sum_s w(v,s) = q(v), \qquad \forall v\in\cV.
    \label{eq:rows-preserved}
\end{align}

\emph{Column sums.}
For each $i$,
\begin{align*}
    \sum_{x\in\cV} \tilde w(x,v_i)
    = &~ \left(w(v_{i-1},v_i)-\varepsilon\right)
       + \left(w(v_i,v_i)+\varepsilon\right)
       + \sum_{x\notin\{v_{i-1},v_i\}} w(x,v_i)\\
    = &~ \sum_{x\in\cV} w(x,v_i),
\end{align*}
where again indices are taken modulo $k$.
All other columns are unchanged, so
\begin{align}
    \sum_v \tilde w(v,s) = \sum_v w(v,s) = p_0(s), \qquad \forall s\in\cV.
    \label{eq:cols-preserved}
\end{align}
Thus $\tilde w\in\cP(p_0,q)$.

\emph{Objective value.}
Only entries on the cycle and the corresponding diagonals change, hence
\begin{align}
    J(\tilde w)-J(w)
    = &~ \sum_{i=0}^{k-1}
       \Bigl[ \tilde w(v_i,v_i)M(v_i,v_i) + \tilde w(v_i,v_{i+1})M(v_i,v_{i+1}) \\
    &~ - w(v_i,v_i)M(v_i,v_i) - w(v_i,v_{i+1})M(v_i,v_{i+1}) \Bigr] \notag\\
    = &~ \varepsilon \sum_{i=0}^{k-1} \left(M(v_i,v_i) - M(v_i,v_{i+1})\right).
    \label{eq:J-diff-cycle}
\end{align}
Let $\sigma\in S_n$ be the permutation whose cycle on the set
$\{v_0,\dots,v_{k-1}\}$ is $(v_0 v_1 \dots v_{k-1})$ and which fixes all
other vertices.
For $v=v_0$ we have $\kappa_v=k$ and
\begin{align*}
    \sigma^i(v_0) = v_i,\quad
    \sigma^{i+1}(v_0)=v_{i+1},\quad i=0,\dots,k-1,
\end{align*}
so by Eq.~\eqref{eq:M-cycle-cond},
\begin{align}
    \sum_{i=0}^{k-1} M(v_i,v_i)
    = &~ \sum_{i=0}^{\kappa_v-1} M(\sigma^i(v_0),\sigma^i(v_0))
    \notag\\
    \geq &~ \sum_{i=0}^{\kappa_v-1} M(\sigma^i(v_0),\sigma^{i+1}(v_0))
     = \sum_{i=0}^{k-1} M(v_i,v_{i+1}).
    \label{eq:cycle-ineq-specialized}
\end{align}
Combining Eq.~\eqref{eq:J-diff-cycle} and Eq.~\eqref{eq:cycle-ineq-specialized} yields
\begin{align}
    J(\tilde w) - J(w)  \ge  0.
    \label{eq:J-nondecrease-cycle}
\end{align}

Now start from the optimizer $w^{(0)}$.  
If $G_{w^{(0)}}$ has no directed cycles, set $\bar w = w^{(0)}$.
Otherwise, choose a directed cycle in $G_{w^{(0)}}$, apply the above
transformation with maximal $\varepsilon$ as chosen above, and obtain
$w^{(1)}\in\cP(p_0,q)$ with $J(w^{(1)})\ge J(w^{(0)})$.
Since $w^{(0)}$ is optimal, $J(w^{(1)})=J(w^{(0)})$, so $w^{(1)}$ is also
optimal.  Moreover, at least one edge of the chosen cycle has
$w^{(1)}(v_i,v_{i+1})=0$.

Iterating this construction, we obtain a sequence of optimal plans
$w^{(m)}$ in $\cP(p_0,q)$ in which the set of off-diagonal edges with
positive mass strictly decreases whenever there is a directed cycle.
Since there are only finitely many off-diagonal entries, this procedure
must terminate.  We thus obtain an optimal plan $\bar w$ such that
$G_{\bar w}$ contains no directed cycle.

\medskip
Define the off-diagonal flow\begin{align*}
    F(v,s) :=
    \begin{cases}
        \bar w(v,s), & v\neq s,\\
        0, & v=s.
    \end{cases}
\end{align*}
For $v\in\cV$ define
\begin{align*}
    \text{out}(v) := \sum_{s\neq v} F(v,s),
    \qquad
    \text{in}(v)  := \sum_{u\neq v} F(u,v).
\end{align*}
Using the marginal constraints of $\bar w$,
\begin{align*}
    q(v) = &~ \sum_s \bar w(v,s)
          = \bar w(v,v) + \text{out}(v),\\
    p_0(v) = &~ \sum_u \bar w(u,v)
            = \bar w(v,v) + \text{in}(v),
\end{align*}
so
\begin{align}
    \text{out}(v) - \text{in}(v)
    = q(v)-p_0(v)
    =
    \begin{cases}
        \frac{\delta}{2}, & v=a,\\ 
        -\frac{\delta}{2}, & v=b,\\
        0, & v\notin\{a,b\}.
    \end{cases}
    \label{eq:flow-balance}
\end{align}
Thus $F$ is a nonnegative flow of value $\delta/2$ from $a$ to $b$ on the
directed acyclic graph $G_{\bar w}$.

We now decompose $F$ into simple $a$–$b$ paths.
Define $F^{(0)}:=F$ and proceed inductively.
Assume $F^{(m)}$ is a nonnegative flow from $a$ to $b$ with balance
equation Eq.~\eqref{eq:flow-balance}.
Since $\text{out}(a)-\text{in}(a)=\delta/2>0$, there exists
$s_1\neq a$ with $F^{(m)}(a,s_1)>0$.
Set $v_0^{(m)}:=a$ and $v_1^{(m)}:=s_1$.


Suppose $v_0^{(m)} = a, \dots, v_j^{(m)}\neq b$ are already recursively constructed such that we use Eq.~\eqref{eq:flow-balance} and nonnegativity to obtain for $j\ge1$,
\begin{align*}
    \text{in}(v_j^{(m)}) \ge F^{(m)}(v_{j-1}^{(m)},v_j^{(m)}) >0.
\end{align*}
For $v_j^{(m)}\notin\{a,b\}$, Eq.~\eqref{eq:flow-balance} gives
$\text{out}(v_j^{(m)}) = \text{in}(v_j^{(m)}) > 0$, so there exists
$v_{j+1}^{(m)}\neq v_j^{(m)}$ with $F^{(m)}(v_j^{(m)},v_{j+1}^{(m)})>0$.
Since $G_{\bar w}$ is acyclic and finite, the sequence
$v_0^{(m)},v_1^{(m)},\dots$ cannot visit a vertex twice; hence the process
must terminate at a vertex with no outgoing edges.  By
\eqref{eq:flow-balance}, such a vertex must satisfy
$\text{out}(v)-\text{in}(v)\le0$.  
Because all intermediate vertices have balance $0$, the only possible
terminal vertex is $b$.
Thus we obtain a simple path
\begin{align*}
    P^{m+1}: a = v_0^{(m)} \to v_1^{(m)} \to \dots \to v_{k_{m}}^{(m)} = b.
\end{align*}

Let
\begin{align*}
    \alpha_{m} := \inf_{0\le i\le k_{m}-1}
                    F^{(m)}(v_i^{(m)},v_{i+1}^{(m)}) >0,
\end{align*}
and define
\begin{align*}
    F^{(m+1)}(v,s)
    := F^{(m)}(v,s)
       - \alpha_{m}\cdot 
         \mathbf{1}\{(v,s)=(v_i^{(m)},v_{i+1}^{(m)}) 
                         \text{ for some }i=0,\dots,k_m-1\}.
\end{align*}
Then $F^{(m+1)}\ge0$ and satisfies the same balance equations Eq.~\eqref{eq:flow-balance}, but has strictly smaller total flow
$\sum_{v,s}F^{(m+1)}(v,s) = \sum_{v,s}F^{(m)}(v,s) - \alpha_{m}$.

Since the total flow is initially $\delta/2$ and decreases by a positive
amount at each step, this procedure terminates after some $L$ steps with
$F^{(L)}\equiv 0$.  
For $\ell=1,\dots,L$, we obtain simple paths
\begin{align*}
    P^\ell: a = v_0^\ell \to v_1^\ell \to \dots \to v_{k_\ell}^\ell=b,
    \qquad \ell=1,\dots,L,
\end{align*}
and coefficients $\alpha_\ell>0$ such that
\begin{align}
    F(v,s)
    = &~ \sum_{\ell=0}^{L-1} \alpha_\ell
       \mathbf{1}\{(v,s)=(v_i^\ell,v_{i+1}^\ell)
                       \text{ for some }i=0,\dots,k_\ell-1\},
       && v\neq s,
       \label{eq:F-decomp}\\
    \sum_{\ell=0}^{L-1} \alpha_\ell = &~ \frac{\delta}{2}.
    \label{eq:alpha-sum}
\end{align}

\medskip

Finally, we prove that the optimal $G_{w^*}$ contains only a single path and derive the closed form solution $w^*$.

For a general feasible $w\in\cP(p_0,q)$, the column constraints give
\begin{align}
    w(v,v)
    = p_0(v) - \sum_{u\neq v} w(u,v), \qquad v\in\cV.
    \label{eq:diag-from-offdiag}
\end{align}
Thus
\begin{align}
    J(w)
    = &~ \sum_{v} w(v,v)M(v,v) + \sum_{u\neq v} w(u,v)M(u,v) \notag\\
    = &~ \sum_v \left(p_0(v)-\sum_{u\neq v}w(u,v)\right)M(v,v)
       +\sum_{u\neq v}w(u,v)M(u,v) \notag\\
    = &~ \underbrace{\sum_v p_0(v)M(v,v)}_{=:J_0}
       + \sum_{u\neq v} w(u,v)\left(M(u,v)-M(v,v)\right).
    \label{eq:J-decomp-general}
\end{align}
For the optimal plan $\bar w$, the off-diagonal part is $F$, so
\begin{align}
    J(\bar w)
    = J_0 + \sum_{u\neq v} F(u,v)\left(M(u,v)-M(v,v)\right).
    \label{eq:J-barp-F}
\end{align}

For each path $P^\ell$, define its \emph{gain}
\begin{align}
    W(P^\ell)
    := \sum_{i=0}^{k_\ell-1}
       \left(
           M(v_i^\ell,v_{i+1}^\ell)
           - M(v_{i+1}^\ell,v_{i+1}^\ell)
       \right).
    \label{eq:path-gain}
\end{align}
Using the decomposition Eq.~\eqref{eq:F-decomp}, we obtain from Eq.~\eqref{eq:J-barp-F}
\begin{align}
    J(\bar w)
    = &~ J_0
       + \sum_{u\neq v}
         \left(
           \sum_{\ell=1}^L \alpha_\ell
           \mathbf{1}\{(u,v)=(v_i^\ell,v_{i+1}^\ell)
                           \text{ for some }i = 0,...,k_\ell-1\}
         \right)
         \left(M(u,v)-M(v,v)\right) \notag\\
    = &~ J_0 + \sum_{\ell=1}^L \alpha_\ell
       \sum_{i=0}^{k_\ell-1}
       \left(
           M(v_i^\ell,v_{i+1}^\ell)
           - M(v_{i+1}^\ell,v_{i+1}^\ell)
       \right) \notag\\
    = &~ J_0 + \sum_{\ell=1}^L \alpha_\ell W(P^\ell).
    \label{eq:J-barp-pathsum}
\end{align}
By Eq.~\eqref{eq:alpha-sum},
\begin{align}
    J(\bar w)
    = J_0 + \frac{\delta}{2} \sum_{\ell=1}^L \theta_\ell W(P^\ell),
    \qquad
    \theta_\ell := \frac{\alpha_\ell}{\delta/2},\quad
    \theta_\ell\ge0,\ \sum_\ell \theta_\ell=1.
    \label{eq:J-barp-convex}
\end{align}

Let
\begin{align}
    W^* := \sup\{W(P): P \text{ is a simple directed path from }a\text{ to }b\}.
    \label{eq:W-star-def}
\end{align}
Since the graph on $\cV$ is finite, the maximum exists.
From Eq.~\eqref{eq:J-barp-convex} we deduce
\begin{align}
    J(\bar w)
    \le J_0 + \frac{\delta}{2} W^*.
    \label{eq:J-barp-upper}
\end{align}


Now fix an arbitrary simple directed path
\begin{align*}
    P: a=u_0 \to u_1 \to \dots \to u_K = b
\end{align*}
with distinct vertices $u_0,\dots,u_K$.
Define $w^P$ by
\begin{align}
    w^P
    := \frac{\delta}{2}\sum_{i=0}^{K-1}
        \left(\mathbf{e}_{u_i}\mathbf{e}_{u_{i+1}}^\top
              - \mathbf{e}_{u_{i+1}}\mathbf{e}_{u_{i+1}}^\top\right)
       + \diag(p_0).
    \label{eq:pP-def}
\end{align}

We first verify $w^P\in\cP(p_0,q)$.

The diagonal entries of $w^P$ are
\begin{align}
    w^P(u_0,u_0) = &~ p_0(u_0), \notag\\
    w^P(u_j,u_j) = &~ p_0(u_j) - \frac{\delta}{2}, \quad 1\le j\le K,
       \label{eq:pP-diag}\\
    w^P(v,v) = &~ p_0(v), \quad v\notin\{u_0,\dots,u_K\}. \notag
\end{align}
Since $\inf_v p_0(v)>\delta$, we have
$p_0(u_j)-\delta/2 > \delta/2>0$ for $1\le j\le K$; thus all diagonals are
nonnegative.
Off-diagonal entries are either $0$ or $\delta/2$, so $w^P\ge0$.

For $v=u_0=a$,
\begin{align*}
    \sum_s w^P(a,s)
    = &~ w^P(a,a) + w^P(a,u_1)
     = p_0(a) + \frac{\delta}{2}
     = q(a).
\end{align*}
For an internal vertex $u_j$ with $1\le j\le K-1$,
\begin{align*}
    \sum_s w^P(u_j,s)
    = &~ w^P(u_j,u_j) + w^P(u_j,u_{j+1})\\
    = &~ \left(p_0(u_j)-\frac{\delta}{2}\right) + \frac{\delta}{2}
     = p_0(u_j) = q(u_j).
\end{align*}
For $v=u_K=b$,
\begin{align*}
    \sum_s w^P(b,s)
    = &~ w^P(b,b)
     = p_0(b)-\frac{\delta}{2}
     = q(b).
\end{align*}
For $v\notin\{u_0,\dots,u_K\}$, the only nonzero entry in row $v$ is the
diagonal, so
\begin{align*}
    \sum_s w^P(v,s) = p_0(v) = q(v).
\end{align*}
Thus the row constraints are satisfied.

For a vertex $u_{i+1}$ on the path (with $0\le i\le K-1$),
\begin{align*}
    \sum_v w^P(v,u_{i+1})
    = &~ w^P(u_{i+1},u_{i+1}) + w^P(u_i,u_{i+1})\\
    = &~ \left(p_0(u_{i+1})-\frac{\delta}{2}\right) + \frac{\delta}{2}
     = p_0(u_{i+1}).
\end{align*}
For $s\notin\{u_1,\dots,u_K\}$, the only nonzero entry in column $s$ is
$w^P(s,s)=p_0(s)$, so
\begin{align*}
    \sum_v w^P(v,s)=p_0(s).
\end{align*}
Hence $w^P\in\cP(p_0,q)$.


From Eq.~\eqref{eq:pP-def} and linearity of $J$,
\begin{align}
    J(w^P) - J_0
    = &~ \frac{\delta}{2}\sum_{i=0}^{K-1}
       \left(
           M(u_i,u_{i+1}) - M(u_{i+1},u_{i+1})
       \right)
     = \frac{\delta}{2} W(P),
    \label{eq:J-pP}
\end{align}
where $W(P)$ is defined as in Eq.~\eqref{eq:path-gain} for this path $P$.

Now choose a path $P^*$ attaining the maximum gain $W(P^*)=W^*$ in Eq.~\eqref{eq:W-star-def}, and set $w^*:=w^{P^*}$.  Then
\begin{align}
    J(w^*)
    = &~ J_0 + \frac{\delta}{2} W^*
     \ge J_0 + \frac{\delta}{2} W(P^\ell)
     \quad\forall \ell,
     \notag\\
    \geq &~ J_0 + \frac{\delta}{2}\sum_{\ell=1}^L \theta_\ell W(P^\ell)
     = J(\bar w)
     \ge J(w),\quad \forall p\in\cP(p_0,q),
    \label{eq:pstar-optimal}
\end{align}
where we used Eq.~\eqref{eq:J-barp-convex} and Eq.~\eqref{eq:J-barp-upper} in the
last line.  
Thus, $w^*$ is an optimizer of $J$ over $\cP(p_0,q)$ with only a single path.


Write the maximizing path as
\begin{align*}
    P^*: a=u_0 \to u_1 \to \dots \to u_K=b.
\end{align*}
Define a permutation $\sigma\in S_n$ by
\begin{align*}
    \sigma(u_i) = &~ u_{i+1}, \quad i=0,\dots,K-1,\\
    \sigma(u_K) = &~ u_0,\\
    \sigma(v) = &~ v, \quad v\notin\{u_0,\dots,u_K\}.
\end{align*}
Then the orbit of $a$ under $\sigma$ is
\begin{align*}
    a, \sigma(a), \dots, \sigma^K(a)
    = u_0,u_1,\dots,u_K,
\end{align*}
and $\sigma^{K+1}(a)=a$, so the minimal positive integer
$\kappa$ with $\sigma^\kappa(a)=a$ is $\kappa=K+1$.
In particular,
\begin{align*}
    \sigma^i(a) = u_i,\quad i=0,\dots,K.
\end{align*}
Therefore, the definition Eq.~\eqref{eq:pP-def} of $w^*$ can be rewritten as
\begin{align*}
    w^*
    = &~ \frac{\delta}{2}\sum_{i=0}^{K-1}
        \left(\mathbf{e}_{u_i}\mathbf{e}_{u_{i+1}}^\top
              - \mathbf{e}_{u_{i+1}}\mathbf{e}_{u_{i+1}}^\top\right)
       + \diag(p_0)\\
    = &~ \frac{\delta}{2}\sum_{i=0}^{\kappa-2}
        \left(\mathbf{e}_{\sigma^i(a)}\mathbf{e}_{\sigma^{i+1}(a)}^\top
              - \mathbf{e}_{\sigma^{i+1}(a)}\mathbf{e}_{\sigma^{i+1}(a)}^\top\right)
       + \diag(p_0),
\end{align*}
which is exactly Eq.~\eqref{eq:p-star-form}.  This completes the proof.
\end{proof}


\begin{lemma}[Optimizer of Middle Problem]\label{lem:middle}
Let $\cV$ be a the set $\{1,\dots,n\}$, $p_0 \in \Delta(\cV)$ be a distribution over $\cV$ such that $\inf_{v \in \cV}p_0(v) > \delta$, and $M: \cV \times \cV \mapsto \R$ be a matrix. Consider the optimization problem
\begin{align*}
    J: \inf_{q \in \cQ(p_0,\delta)} \sup_{w \in \cP(p_0,q)} \sum_{v,s \in \cV} w(v,s) \cdot M(v,s)
\end{align*}
where 
\begin{align*}
    \cQ(p_0,\delta) = &~ \left\{q \in \Delta(\cV): \|q - p_0\|_1 \leq \delta\right\} \\
    \cP(p_0,q) = &~ \left\{w(v,s): \sum_{s}w(v,s) = q(v), \sum_vw(v,s) =  p_0(s)\right\}
\end{align*}
Then there exists an optimal solution $q^*$ of $J$ coming from the set 
\begin{align*}
    \cQ_{\ext}(p_0,\delta) := \left\{p_0 + \frac{\delta}{2}\cdot (\mathbf{e}_a - \mathbf{e}_b): a \neq b \in \cV\right\}.
\end{align*}
\end{lemma}

\begin{proof}
By \Cref{lem:q_concave_vertices}, $\cQ_{\ext}(p_0,\delta)$ is the set of extreme points of the convex polytope $\cQ(p_0,\delta)$. 
Define $J(q) = \sup_{w \in \cP(p_0,q)} \sum_{v,s \in \cV} w(v,s) \cdot M(v,s)$. We show that $J(q)$ is a concave function. Indeed for all $q_1 \neq q_2$, we let 
\begin{align*}
    w_1 := &~ \arg \sup_{w \in \cP(p_0,q_1)} \sum_{v,s \in \cV} w(v,s) \cdot M(v,s)\\
    w_2 := &~ \arg \sup_{w \in \cP(p_0,q_2)} \sum_{v,s \in \cV} w(v,s) \cdot M(v,s).
\end{align*}
We have
\begin{align*}
    J((q_1+q_2)/2) = &~ \sup_{w \in \cP(p_0,(q_1+q_2)/2)} \sum_{v,s \in \cV} w(v,s) \cdot M(v,s) \\
    \geq &~ \sum_{v,s \in \cV} (w_1(v,s) + w_2(v,s)) /2 \cdot M(v,s)\\
    = &~ (J(q_1)+J(q_2))/2,
\end{align*}
where the second step is because the feasibility of $(w_1(v,s) + w_2(v,s))/2$:
\begin{align*}
    \sum_{s}(w_1(v,s) + w_2(v,s))/2 = (q_1(v)+q_2(v))/2.
\end{align*}
Since $J(q)$ is concave over a convex feasible set, its optimizer must be found on the extreme set $\cQ_{\ext}(p_0,\delta)$. This completes the proof.
\end{proof}


\begin{lemma}[Optimal Subpaths]\label{lem:optimalpaths}
Let $\cV$ be a the set $\{1,\dots,n\}$, $p_0 \in \Delta(\cV)$ be a distribution over $\cV$ such that $\inf_{v \in \cV}p_0(v) > \delta$.
Let $M: \cV \times \cV \mapsto \R$ be a matrix that satisfies
\begin{align*}
    \sum_{i=0}^{\kappa_v-1} M(\sigma^i(v),\sigma^i(v)) 
     \geq  \sum_{i=0}^{\kappa_v-1} M(\sigma^i(v),\sigma^{i+1}(v))
\end{align*}
for any $v \in \cV$, $\sigma \in S_n$, where $\kappa_v = \inf\{k\ge1:\sigma^{k}(v) = v\}$. 
Consider the optimization problem
$$J(q) = \sup_{w \in \cP(p_0,q)}\sum_{v,s \in \cV}w(v,s)\cdot M(v,s)$$
where 
$$\cP(p_0,q) = \left\{w(v,s): \sum_{s}w(v,s) = q(v), \sum_vw(v,s) =  p_0(s)\right\}.$$
If there exists $q = p_0 + \frac{\delta}{2}\cdot(\mathbf{e}_a-\mathbf{e}_b)$ for some $a \neq b \in \cV$ and $w^*_q$ such that $w^*_q$ is the optimizer of $J(q)$ and for a permutation $\sigma$ we have that
$$w^*(q) = \frac{\delta}{2} \cdot \sum_{i=0}^{\kappa-2} (\mathbf{e}_{\sigma^{i}(a)}\mathbf{e}_{\sigma^{i+1}(a)}^\top - \mathbf{e}_{\sigma^{i+1}(a)}\mathbf{e}_{\sigma^{i+1}(a)}^\top) + \diag(p_0),$$
where $\kappa$ is the minimum positive integer such that $\sigma^{\kappa}(a) = a$. Then define $c := \sigma^{i}(a)$ and $d := \sigma^{i+1}(a)$ with $i < \kappa$, and let $q' = p_0 + \frac{\delta}{2}\cdot (\mathbf{e}_c - \mathbf{e}_d)$. We have that the $w^*(q')$ that optimizes $J(q')$ is given by
$$w^*(q') = \frac{\delta}{2}\cdot(\mathbf{e}_{c}\mathbf{e}_d^\top - \mathbf{e}_d \mathbf{e}_d^\top) + \diag(p_0).$$
\end{lemma}

\begin{proof}
Let $c = \sigma^i(a)$ and $d = \sigma^{i+1}(a)$. Define the local transport term corresponding to the edge $(c, d)$ as:
$$ T_{c,d} := \frac{\delta}{2} \cdot (\mathbf{e}_c \mathbf{e}_d^\top - \mathbf{e}_d \mathbf{e}_d^\top). $$
Note that the proposed optimizer for the subproblem $J(q')$ is given by $\hat{w} = \diag(p_0) + T_{c,d}$. 

We proceed by contradiction. Assume that $\hat{w}$ is not the optimizer for $J(q')$. Then there exists a feasible transport plan $\tilde{w} \in \cP(p_0, q')$ such that:
$$ \sum_{v,s \in \cV} \tilde{w}(v,s) M(v,s) > \sum_{v,s \in \cV} \hat{w}(v,s) M(v,s). $$
By \Cref{lem:inner}, $\tilde{w}$ can be chosen to take the form of
\begin{align*}
    \tilde{w} = \frac{\delta}{2} \sum_{i=0}^{\bar \kappa-2} 
        \left(\mathbf{e}_{\bar \sigma^{i}(c)}\mathbf{e}_{\sigma^{i+1}(c)}^\top 
              - \mathbf{e}_{\bar \sigma^{i+1}(c)}\mathbf{e}_{\bar \sigma^{i+1}(c)}^\top\right) 
      + \diag(p_0)
\end{align*}
for some $\bar \sigma \in S_n$ and $\bar \kappa = \inf\{k:\bar \sigma^k(c) = c\}$.
It follows that $\inf_{v \in \cV}\left(\tilde{w} - \diag(p_0)\right)(v,v) \geq -\delta/2$ and $\inf_{v \neq s \in \cV}\left(\tilde{w} - \diag(p_0)\right)(v,s) \geq 0$.

Substituting $\hat{w} = \diag(p_0) + T_{c,d}$ into the inequality, we have:
\begin{equation} \label{eq:contra_ineq}
\sum_{v,s \in \cV} (\tilde{w}(v,s) - \diag(p_0)_{v,s}) M(v,s) > \sum_{v,s \in \cV} (T_{c,d})_{v,s} M(v,s).
\end{equation}

Now, consider the global optimizer $w^*(q)$. By the hypothesis, $w^*(q)$ decomposes into a sum of path segments. We can separate the specific term $T_{c,d}$ from the rest of the path:
$$ w^*(q) = \underbrace{\left( \diag(p_0) + \frac{\delta}{2} \sum_{j \neq i}^{\kappa-2} (\mathbf{e}_{\sigma^{j}(a)}\mathbf{e}_{\sigma^{j+1}(a)}^\top - \mathbf{e}_{\sigma^{j+1}(a)}\mathbf{e}_{\sigma^{j+1}(a)}^\top) \right)}_{R} + T_{c,d}, $$
where $R$ represents the flow on the path excluding the step from $d$ to $c$. We construct a new global transport plan $w_{\text{new}}$ by replacing the local step $T_{c,d}$ in $w^*(q)$ with the ``better'' local flow derived from $\tilde{w}$:
$$ w_{\text{new}} := R + (\tilde{w} - \diag(p_0)). $$
Substituting $R = w^*(q) - T_{c,d}$, we get:
$$ w_{\text{new}} = w^*(q) - T_{c,d} + \tilde{w} - \diag(p_0). $$

We verify that $w_{\text{new}}$ is feasible for the original problem $J(q)$:

Since $\tilde{w} \in \cP(p_0, q')$, its row sum is $q' = p_0 + \frac{\delta}{2}(\mathbf{e}_c - \mathbf{e}_d)$. The term $T_{c,d}$ also corresponds to a row marginal shift of $\frac{\delta}{2}(\mathbf{e}_c - \mathbf{e}_d)$. Thus, $w_{\text{new}}$ preserves the row sums of $w^*(q)$, which equal $q$.
 
Both $\tilde{w}$ and $\diag(p_0) + T_{c,d}$ maintain column sums equal to $p_0$. Thus, $w_{\text{new}}$ maintains the column sums of $w^*(q)$, which equal $p_0$.
 
Since $\inf_{v} p_0(v) > \delta$, $\inf_{v} T_{c,d}(v,v) \geq -\delta/2$ , and $\inf_{v} \left(\tilde{w} - \diag(p_0)\right)(v,v) \geq -\delta/2$, we have $\inf_{v} w_{\text{new}}(v,v) \geq 0$. Furthermore, all off-diagonal entries of $w^*(q) - T_{c,d}$ and $\tilde{w} - \diag(p_0)$ are non-negative, thus we conclude that $w_{\text{new}}$ is non-negative.


Finally, we compare the objective value of $w_{\text{new}}$ to $w^*(q)$:
\begin{align*}
J(w_{\text{new}}) = &~ \sum_{v,s} (R_{v,s} + \tilde{w}(v,s) - \diag(p_0)_{v,s}) M(v,s) \\
= &~ \sum_{v,s} R_{v,s} M(v,s) + \sum_{v,s} (\tilde{w}(v,s) - \diag(p_0)_{v,s}) M(v,s) \\
&> \sum_{v,s} R_{v,s} M(v,s) + \sum_{v,s} (T_{c,d})_{v,s} M(v,s) \quad \text{(by Ineq. \ref{eq:contra_ineq})} \\
= &~ J(w^*(q)).
\end{align*}
We have constructed a feasible solution $w_{\text{new}} \in \cP(p_0, q)$ with a strictly higher objective value than $w^*(q)$. This contradicts the optimality of $w^*(q)$. Therefore, $\hat{w}$ must be the optimizer for $J(q')$.
\end{proof}

\begin{corollary}[Equivalence]\label{cor:equiv}
Let $\cV$ be a the set $\{1,\dots,n\}$ and $p_0 \in \Delta(\cV)$ be a distribution over $\cV$ such that $\inf_{v \in \cV}p_0(v) > \delta$. The following problem
\begin{align*}
    \sup_{r} \inf_{q \in \cQ(p_0,\delta)} \sup_{w \in \cP(p_0,q)} &~  \sum_{v,s \in \cV} w(v,s) \cdot \log r(v,s) \\
    s.t. &~ \sum_{s \in \cV} r(v,s) = 1, ~ \forall v \in \cV\\
    &~ r(v,s) \geq 0,~ \forall v,s \in \cV \notag 
\end{align*}
where 
\begin{align*}
    \cQ(p_0,\delta) = &~ \left\{q \in \Delta(\cV): \|q - p_0\|_1 \leq \delta\right\} \\
    \cP(p_0,q) = &~ \left\{w(v,s): \sum_{s \in \cV}w(v,s) = q(v), \sum_{v  \in \cV}w(v,s) =  p_0(s)\right\}
\end{align*}
is equivalent to the following problem
\begin{align*}
    \sup_{r} \inf_{q \in \cQ_{\text{ext}}(p_0,\delta)} \sup_{w \in \cP_{\text{flow}}(p_0,q)} &~  \sum_{v,s \in \cV} w(v,s) \cdot \log r(v,s) \\
    s.t. &~ \sum_{s \in \cV} r(v,s) = 1, ~ \forall v \in \cV\\
    &~ \sum_{i=0}^{n-1} \log r(v,s)(\sigma^i(v),\sigma^i(v)) \geq \sum_{i=0}^{n-1} \log r(v,s)(\sigma^i(v),\sigma^{i+1}(v)), ~\forall \sigma  \in S_n,\\
    &~ r(v,s) \geq 0,~ \forall v,s \in \cV \notag 
\end{align*}
where 
\begin{align*}
    \cQ_{\ext}(p_0,\delta) = &~ \left\{p_0 + \frac{\delta}{2}\cdot (\mathbf{e}_a - \mathbf{e}_b): a \neq b \in \cV\right\} \\
    \cP_{\flow}(p_0,q) = &~ \left\{\frac{\delta}{2} \cdot \sum_{i=0}^{\kappa-2} (\mathbf{e}_{\sigma^{i}(a)}\mathbf{e}_{\sigma^{i+1}(a)}^\top - \mathbf{e}_{\sigma^{i+1}(a)}\mathbf{e}_{\sigma^{i+1}(a)}^\top) + \diag(p_0): a,b \in \cV, \kappa = \inf \{k \in \Z_+: \sigma^k(a) = a\}\right\}
\end{align*}
\end{corollary}

\begin{proof}
Let 
\begin{align*}
    \cR = &~ \left\{r \in \R^{n \times n}: \sum_{s \in \cV} r(v,s) = 1, ~r(v,s) \geq 0,~ \forall v,s \in \cV \right\}\\
    \cR_{\diag} = &~ \cR \cap \left\{r:\sum_{i=0}^{n-1} \log r(v,s)(\sigma^i(v),\sigma^i(v)) \geq \sum_{i=0}^{n-1} \log r(v,s)(\sigma^i(v),\sigma^{i+1}(v)), ~\forall \sigma \in S_n \right\}.
\end{align*}
Combining \Cref{lem:diagonal-dom}, \Cref{lem:inner}, and \Cref{lem:middle}, we have
\begin{align*}
    \sup_{r \in \cR} \inf_{q \in \cQ(p_0,\delta)} \sup_{w \in \cP(p_0,q)} \sum_{v,s \in \cV} w(v,s) \cdot \log r(v,s) = &~ \sup_{r \in \cR_{\diag}} \inf_{q \in \cQ(p_0,\delta)} \sup_{w \in \cP(p_0,q)} \sum_{v,s \in \cV} w(v,s) \cdot \log r(v,s)\\
    = &~ \sup_{r \in \cR} \inf_{q \in \cQ_{\ext}(p_0,\delta)} \sup_{w \in \cP(p_0,q)} \sum_{v,s \in \cV} w(v,s) \cdot \log r(v,s)\\
    = &~ \sup_{r \in \cR} \inf_{q \in \cQ_{\ext}(p_0,\delta)} \sup_{w \in \cP_{\flow}(p_0,q)} \sum_{v,s \in \cV} w(v,s) \cdot \log r(v,s)
\end{align*}
where the first step uses \Cref{lem:diagonal-dom} and the fact that permuting $\cV$ and the columns of $r$ by the same $\sigma \in S_n$ does not change the objective, the second step uses \Cref{lem:middle}, and the last step uses \Cref{lem:inner}. This completes the proof.
\end{proof}

\begin{proposition}[Reformulated Problem]\label{prop:reform-problem}
Let $\cV$ be a the set $\{1,\dots,n\}$ and $p_0 \in \Delta(\cV)$ be a distribution over $\cV$ such that $\inf_{v \in \cV}p_0(v) > \delta$. Consider the following problem
\begin{align}\label{eq:gen_opt_reform}
    \sup_{r} \inf_{q \in \cQ(p_0,\delta)} \sup_{w \in \cP(p_0,q)} &~  \sum_{v,s \in \cV} w(v,s) \cdot \log r(v,s) \\
    s.t. &~ \sum_{s \in \cV} r(v,s) = 1, ~ \forall v \in \cV\\
    &~ r(v,s) \geq 0,~ \forall v,s \in \cV \notag 
\end{align}
where $p_0$ is a fixed distribution over the sample space $\cV$ such that $\inf_{v \in \cV}p_0(v) > \delta$, and
\begin{align*}
    \cQ(p_0,\delta) = &~ \left\{q \in \Delta(\cV): \|q - p_0\|_1 \leq \delta\right\} \\
    \cP(p_0,q) = &~ \left\{w(v,s): \sum_{s \in \cV}w(v,s) = q(v), \sum_{v  \in \cV}w(v,s) =  p_0(s)\right\}.
\end{align*}
Then the optimum value is
\begin{align*}
J^*
= (1-\tfrac{\delta}{2})\log\left(1-\frac{\delta}{2}\right)
 + \tfrac{\delta}{2}\log\left(\frac{\delta}{2(n-1)}\right).
\end{align*}
In particular, this is achieved at
\begin{align*}
r^*_{a,b}(v,s)
=
\begin{cases}
1-\delta/2, & s=v,\\
\delta/(2n-2), & s\neq v,
\end{cases}
\end{align*}
and for any $q \in \cQ(p_0,\delta)$ written as $q = \sum_{i=1}^k \lambda_i \cdot q_i$ for 
$$q_i = p_0 + \frac{\delta}{2}\cdot (\mathbf{e}_{v_i} - \mathbf{e}_{s_i}) \in \cQ_{\ext}(p_0,\delta) := \left\{p_0 + \frac{\delta}{2}\cdot (\mathbf{e}_a - \mathbf{e}_b): a \neq b \in \cV\right\}$$
the optimizer of the inner problem is given by
\begin{align*}
    \sum_{i=1}^k \lambda_i \cdot \left(\frac{\delta}{2} \cdot (\mathbf{e}_{v_i}\mathbf{e}_{s_i}^\top - \mathbf{e}_{s_i}\mathbf{e}_{s_i}^\top) + \diag(p_0)\right).
\end{align*}
\end{proposition}

\begin{proof}

By \Cref{cor:equiv}, WLOG it suffices to consider $q$ in the form of $p_0 + \frac{\delta}{2}\cdot (\mathbf{e}_a - \mathbf{e}_b)$ for $a \neq b \in \cV$ and $p$ written as $\frac{\delta}{2} \cdot \sum_{i=0}^{\kappa-2} (\mathbf{e}_{\sigma^{i}(a)}\mathbf{e}_{\sigma^{i+1}(a)}^\top - \mathbf{e}_{\sigma^{i+1}(a)}\mathbf{e}_{\sigma^{i+1}(a)}^\top) + \diag(p_0)$ for $\sigma \in S_n, \kappa = \inf \{k \in \Z_+: \sigma^k(a) = a\}$ and $a,b \in \cV$.

Define
\begin{align*}
    J(r,w,q) = &~ \sum_{v,s \in \cV} w(v,s) \cdot \log r(v,s)\\
    J^* = &~ \sup_{r} \inf_{q \in \cQ_{\text{ext}}(p_0,\delta)} \sup_{w \in \cP_{\flow}(p_0,q)} J(r,w,q)\\
    R^* = &~ \{r: \inf_{q \in \cQ_{\text{ext}}(p_0,\delta)} \sup_{w \in \cP_{\flow}(p_0,q)} J(r,w,q) = J^*\}\\
    r^* = &~ \arg\sup\left\{\tr[r]: r \in R^* \right\}.
\end{align*}
We say a solution $(\bar w,\bar q)$ is active if 
\begin{align*}
    \bar w = &~ \arg\sup_{w \in \cP_{\flow}(p_0,\bar q)} J(r^*, w,\bar q)\\
    \bar q = &~ \arg\inf_{q \in \cQ_{\text{ext}}(p_0,\delta)} \sup_{w \in \cP_{\flow}(p_0,q)} J(r^*,w,q)\\
    J(\bar w,\bar q) = &~ J^*.
\end{align*}

Fix $v \neq s \in \cV$ and consider the perturbation for sufficiently small $\epsilon > 0$
\begin{align*}
    \bar r(v,v) \leftarrow &~  r^*(v,v) + \epsilon\\
    \bar r(v,s) \leftarrow &~  r^*(v,s) - \epsilon\\
    \bar r(v',s') \leftarrow &~  r^*(v',s'), ~\forall (v',s') \neq (v,s).
\end{align*}
Then $\bar{r}$ cannot be a valid solution since it has higher trace than $r^*$. It follows that
\begin{align*}
    \frac{d J(\bar r,\bar w,\bar q)}{d\epsilon} = \frac{\bar w(v,v)}{r^*(v,v)} - \frac{\bar w(v,s)}{r^*(v,s)} \leq 0, ~\forall \text{ active } (\bar w,\bar q).
\end{align*}
Notice that the derivative $\frac{d J(\bar r, \bar w,\bar q)}{d\epsilon}$ must be either 
\begin{itemize}
    \item No-transport: $\bar w(v,v) = p_0(v), \bar w(v,s) = 0$ so $\frac{d J(\bar r, \bar w,\bar q)}{d\epsilon} = \frac{p_0(v)}{r^*(v,v)} > 0$.
    \item Middle-way: $\bar w(v,v) = p_0(v) - \delta/2, \bar w(v,s) = \delta/2$ so $\frac{d J(\bar r, \bar w,\bar q)}{d\epsilon} = \frac{p_0(v) - \delta/2}{r^*(v,v)} - \frac{\delta/2}{r^*(v,s)}$.
    \item Transport-start: $\bar w(v,v) = p_0(v), \bar w(v,s) = \delta/2$ so $\frac{d J(\bar r, \bar w,\bar q)}{d\epsilon} = \frac{p_0(v)}{r^*(v,v)} - \frac{\delta/2}{r^*(v,s)}$.
    \item Transport-end: $\bar w(v,v) = p_0(v) - \delta/2, \bar w(v,s) = 0$ so $\frac{d J(\bar r, \bar w,\bar q)}{d\epsilon} = \frac{p_0(v) - \delta/2}{r^*(v,v)} > 0$.
\end{itemize}
Since $\frac{d J(\bar r, \bar w,q)}{d\epsilon} < 0$, we can rule out the first and last cases. Now we can summarize that for any $v \neq s \in \cV$ we have either Middle-way: $\bar w(v,v) = p_0(v) - \delta/2, \bar w(v,s) = \delta/2$ or Transport-start: $\bar w(v,v) = p_0(v), \bar w(v,s) = \delta/2$, and in any case $\frac{p_0(v) - \delta/2}{r^*(v,v)}  - \frac{\delta/2}{r^*(v,s)} \leq 0$ holds.

Since $p_0(v) > \delta$, in either Middle-way or Transport-start case we have
\begin{align*}
    0 \geq \frac{p_0(v) - \delta/2}{r^*(v,v)} - \frac{\delta/2}{r^*(v,s)} > \frac{\delta/2}{r^*(v,v)} - \frac{\delta/2}{r^*(v,s)}.
\end{align*}
It follows that $r^*(v,v) > r^*(v,s)$. Since this argument holds for all $v \neq s \in \cV$, $r^*$ must satisfy $r^*(v,v) > r^*(v,s)$ for all $v \neq s \in \cV$. 

Next, we show that all active $\bar q, \bar w$ can be written as $\bar q = p_0 + \frac{\delta}{2}\cdot (\mathbf{e}_v - \mathbf{e}_s), \bar w = \frac{\delta}{2} \cdot (\mathbf{e}_{v}\mathbf{e}_{s}^\top - \mathbf{e}_{s}\mathbf{e}_{s}^\top) + \diag(p_0)$ for some $v\neq s \in \cV$.

In either Middle-way or Transport-start case, there exists $a \neq b \in \cV$ and $\sigma \in S_n$ and $i \in \Z$ such that $\bar q = p_0 + \frac{\delta}{2}\cdot (\mathbf{e}_a - \mathbf{e}_b)$ and $v = \sigma^i(a), s = \sigma^{i+1}(a)$, due to \Cref{lem:inner}. By \Cref{lem:optimalpaths}, with respect to $\bar q = p_0 + \frac{\delta}{2}\cdot(\mathbf{e}_v-\mathbf{e}_s)$, the optimizer of the inner problem 
must be written as 
\begin{align*}
    \frac{\delta}{2}\cdot(\mathbf{e}_v\mathbf{e}_s^\top - \mathbf{e}_s \mathbf{e}_s^\top) + \text{diag}(p_0) = \arg\sup_{w \in \cP_{\flow}(p_0,\bar q)} J(r^*, w,\bar q).
\end{align*}
Since this argument holds for all $v \neq s \in \cV$, we establish a one-to-one correspondence between $\bar q = p_0 + \frac{\delta}{2}\cdot (\mathbf{e}_v - \mathbf{e}_s)$ and $\bar w = \frac{\delta}{2} \cdot (\mathbf{e}_{v}\mathbf{e}_{s}^\top - \mathbf{e}_{s}\mathbf{e}_{s}^\top) + \diag(p_0)$ for all active $\bar q, \bar w$.

We can now explicitly write:
\begin{align*}
    \inf_{q \in \cQ_{\text{ext}}(p_0,\delta)} \sup_{w \in \cP_{\flow}(p_0,q)} J(r^*,w,q) = &~ \inf_{q \in \cQ_{\text{ext}}(p_0,\delta)} \sup_{w \in \cP_{\flow}(p_0,q)} \sum_{v,s \in \cV} w(v,s) \cdot \log r^*(v,s)\\
    = &~ \inf_{v \neq s \in \cV}\sum_{x \in \cV} p_0(x) \log r^*(x,x) - \frac{\delta}{2}\cdot(\log r^*(s,s) - \log r^*(v,s))
\end{align*}
Define $J(r, v,s) = \sum_{x \in \cV} p_0(x) \log r(x,x) - \frac{\delta}{2}\cdot(\log r(s,s) - \log r(v,s))$, we claim that $J(r^*, v,s)$ must be the same for all $v \neq s \in \cV$.
Otherwise suppose
\begin{align*}
    J(r^*, \bar v, \bar s) = \sup_{v,s}J(r^*, v,s) > J^*.
\end{align*}
Set
\begin{align*}
    \bar r(\bar v, \bar v) \leftarrow &~  r^*(\bar v, \bar v) + \epsilon\\
    \bar r(\bar v, \bar s) \leftarrow &~  r^*(\bar v, \bar s) - \epsilon\\
    \bar r(v',s') \leftarrow &~  r^*(v',s'), ~\forall (v',s') \neq (\bar v, \bar s)
\end{align*}
for sufficiently small $\epsilon > 0$. Then for any $(v',s') \neq (\bar v, \bar s)$ we have 
\begin{align*}
    \frac{d J(\bar r, v,s)}{d\epsilon} = \frac{p_0(s) - \delta/2}{r^*(s,s)} \geq 0
\end{align*}
and thus
\begin{align*}
    \inf_{q \in \cQ_{\ext}(p_0,\delta)} \sup_{w \in \cP_{\flow}(p_0,q)} J(r^*,w,q) = \inf_{v \neq s \in \cV, (v,s) \neq (\bar v, \bar s)}J(r^*, v,s) \geq J^*.
\end{align*}
But $\tr[\bar r] > \tr [r^*]$, this is a contradiction.

From the last argument, we know that the optimal solution $r^*$ is of the form
\begin{align*}
r^*_{a,b}(v,s)
=
\begin{cases}
a, & s=v,\\
b, & s\neq v,
\end{cases}
\qquad a\ge 0,\ b\ge 0,\quad a+(n-1)b=1,
\end{align*}
with $a>b$ (strict diagonal dominance). We now optimize over the two parameters $(a,b)$ subject to this constraint:
\begin{align*}
\sup_{a,b}~\Phi(a,b)
\quad\text{s.t. } a+(n-1)b=1,\ a>0,\ b>0,\ a>b.
\end{align*}
straightforward algebra shows the optimal objective value:
\begin{align*}
\Phi(a^*,b^*)
= &~ (1-\tfrac{\delta}{2})\log a^* + \tfrac{\delta}{2}\log b^*\\
= &~ (1-\tfrac{\delta}{2})\log\left(1-\frac{\delta}{2}\right)
 + \tfrac{\delta}{2}\log\left(\frac{\delta}{2(n-1)}\right).
\end{align*}
attained at 
$a^* = 1 - \frac{\delta}{2}$ and $b^* = \frac{\delta}{2(n-1)}.
$
Thus, the optimal value is
\begin{align*}
J^*
= (1-\tfrac{\delta}{2})\log\left(1-\frac{\delta}{2}\right)
 + \tfrac{\delta}{2}\log\left(\frac{\delta}{2(n-1)}\right).
\end{align*}
This completes the proof.

\end{proof}

\begin{lemma}[Row Normalization]\label{lem:normalizationconstant}
Let $\cV$ be a the set $\{1,\dots,n\}$ and $p_0 \in \Delta(\cV)$ be a distribution over $\cV$ such that $\inf_{v \in \cV}p_0(v) > \delta$. Consider the problem
\begin{align*}
    \sup_{e} \inf_{q \in \cQ(p_0,\delta)} \sup_{w \in \cP(p_0,q)} &~  \sum_{v,s \in \cV} w(v,s) \cdot \log e(v,s) \\
    s.t. &~ \sum_{v,s \in \cV} q(v) p_0(s) e(v,s) \leq 1, ~ \forall q \in \cQ\notag\\
    &~ e(v,s) \geq 0,~ \forall v,s \in \cV \notag 
\end{align*}
where 
\begin{align*}
    \cQ(p_0,\delta) = &~ \left\{q \in \Delta(\cV): \|q - p_0\|_1 \leq \delta\right\} \\
    \cP(p_0,q) = &~ \left\{w(v,s): \sum_{s}w(v,s) = q(v), \sum_vw(v,s) =  p_0(s)\right\}
\end{align*}
Now define the kernel matrix $r:\cV \times \cV \mapsto \mathbb{R}$ such that each of its entries are defined as 
$$r(v,s) = \frac{p_0(s)e(v,s)}{A(v)}, \quad \forall v,s \in \cV,$$
where $A(v) := \sum_s p_0(s)e(v,s)$. Then $A^*(v) = \sum_sp_0(s)e^*(v,s) = 1$ at the optimizer $e^*$.
\end{lemma}

\begin{proof}
Let $e(v,s)$ be any feasible solution to the optimization problem. We define the scaling factor of $e$ at node $v$ as:
$$A(v) := \sum_{s \in \cV} p_0(s) e(v,s).$$
Since we are maximizing an objective involving $\log e(v,s)$, we can assume $e(v,s) > 0$ strictly (otherwise the objective is $-\infty$), and consequently $A(v) > 0$.

We can decompose the matrix $e(v,s)$ into a scale-independent "shape" matrix $\bar{e}(v,s)$ and the scaling factors $A(v)$ as follows:
$$e(v,s) = A(v) \cdot \bar{e}(v,s), \quad \text{where } \bar{e}(v,s) = \frac{e(v,s)}{A(v)}.$$
By construction, the normalized matrix $\bar{e}$ satisfies the normalization property:
$$\sum_{s \in \cV} p_0(s) \bar{e}(v,s) = \sum_{s \in \cV} p_0(s) \frac{e(v,s)}{A(v)} = \frac{1}{A(v)} \sum_{s \in \cV} p_0(s) e(v,s) = 1.$$

Now, let us analyze the constraint given in the problem statement. The condition is:
$$\sum_{v,s \in \cV} q(v) p_0(s) e(v,s) \leq 1, \quad \forall q \in \cQ.$$
Substituting the decomposition of $e(v,s)$:
$$\sum_{v \in \cV} q(v) \left( \sum_{s \in \cV} p_0(s) e(v,s) \right) = \sum_{v \in \cV} q(v) A(v) \leq 1, \quad \forall q \in \cQ.$$

Next, we substitute the decomposition into the objective function. Using the property that $w \in \cP(p_0, q)$ implies $\sum_s w(v,s) = q(v)$, we have:
\begin{align*}
\sum_{v,s \in \cV} w(v,s) \log e(v,s) = &~ \sum_{v,s \in \cV} w(v,s) \log \left( A(v) \cdot \bar{e}(v,s) \right) \\
= &~ \sum_{v,s \in \cV} w(v,s) \log \bar{e}(v,s) + \sum_{v,s \in \cV} w(v,s) \log A(v) \\
= &~ \sum_{v,s \in \cV} w(v,s) \log \bar{e}(v,s) + \sum_{v \in \cV} q(v) \log A(v).
\end{align*}
The first term depends only on the normalized shape $\bar{e}$, while the second term depends only on the scaling factors $A(v)$. To maximize the total objective, we must maximize the second term subject to the feasibility constraint derived above.

Consider the term $\sum_{v \in \cV} q(v) \log A(v)$. Since the logarithm is a concave function, we can apply Jensen's inequality:
$$ \sum_{v \in \cV} q(v) \log A(v) \leq \log \left( \sum_{v \in \cV} q(v) A(v) \right). $$
From the feasibility constraint, we know that $\sum_{v \in \cV} q(v) A(v) \leq 1$. Therefore:
$$ \sum_{v \in \cV} q(v) \log A(v) \leq \log(1) = 0. $$

Thus for every $q$,
\begin{align*}
    \sup_{w \in \cP(p_0,q)} \sum_{v,s} w(v,s)\log e(v,s)
     \le 
    \sup_{w \in \cP(p_0,q)} \sum_{v,s} w(v,s)\log \bar e(v,s).
\end{align*}

And the inequality is strict unless $\sum_v q(v)A(v) = 1$ and $A(v)$ is constant on the support of $q$.

Taking the {minimum over $q$}, we conclude:
\begin{align*}
    \inf_q  \sup_p \sum_{v,s} w(v,s)\log e(v,s)
     \le 
    \inf_q  \sup_p \sum_{v,s} w(v,s)\log \bar e(v,s).
\end{align*}

Equality is achieved if and only if $A(v) = 1$ for all $v \in \cV$.
Thus, for any optimal solution $e^*$, the scaling factors must be set to $1$. Consequently:
$$A^*(v) = \sum_{s \in \cV} p_0(s) e^*(v,s) = 1.$$
\end{proof}
